\documentclass[a4paper,11pt]{ltjsarticle}
\usepackage[margin=25mm]{geometry}
\usepackage[backend=biber,style=numeric,sorting=none]{biblatex}
\addbibresource{references.bib}
\usepackage[colorlinks=true,allcolors=blue]{hyperref}

\title{文献管理のサンプル\\[0.4em]\large LuaLaTeX + biblatex + Biber}
\author{サンプル作成者}
\date{2026年10月18日}
\begin{document}
\maketitle
\section{本文から文献を引用する}
biblatex では，文献データを \texttt{references.bib} にまとめ，
本文から \verb|\cite| などを使って引用する\cite{biblatex}。
日本語を含む文献も UTF-8 のまま管理できる。
例えば，仮想の日本語文献を引用すると，次のようになる\cite{sample-japanese}。

\section{Biberを自動実行する}
この例は LuaLaTeX を使う。
latexmk が文献データの状態を見て，Biber と LuaLaTeX を必要な回数だけ実行する。
エディター側で「LuaLaTeX，Biber，LuaLaTeXを2回」のような固定の手順を組む必要はない。

\section{文献ファイルの最小例}
日本語文献のデータは，説明のために作成した架空のものである。
実際のレポートでは，参照した資料の著者・題名・発行年に置き換える。

\printbibliography[title={参考文献}]
\end{document}
